% 公式编写示例
% 本文件展示如何在论文中编写各种数学公式

\chapter{公式编写示例}

\section{行内公式}

行内公式使用 \texttt{\$...\$} 包围，例如：爱因斯坦质能方程 $E=mc^2$，欧拉公式 $e^{i\pi}+1=0$。

常用的行内公式：
\begin{itemize}
  \item 变量：$x$, $y$, $z$
  \item 上标：$x^2$, $x^{n+1}$
  \item 下标：$x_1$, $x_{i,j}$
  \item 分数：$\frac{1}{2}$, $\frac{a}{b}$
  \item 根号：$\sqrt{2}$, $\sqrt[3]{8}$
\end{itemize}

\section{行间公式}

\subsection{无编号公式}

使用 \texttt{\textbackslash[...\textbackslash]} 或 \texttt{equation*} 环境：

\[
  f(x) = ax^2 + bx + c
\]

\begin{equation*}
  \int_0^1 x^2 \, dx = \frac{1}{3}
\end{equation*}

\subsection{带编号公式}

使用 \texttt{equation} 环境：

\begin{equation}
  E = mc^2
  \label{eq:einstein}
\end{equation}

可以引用公式 \eqref{eq:einstein}。

\section{常用数学符号}

\subsection{希腊字母}

小写：$\alpha$, $\beta$, $\gamma$, $\delta$, $\epsilon$, $\theta$, $\lambda$, $\mu$, $\pi$, $\sigma$, $\omega$

大写：$\Gamma$, $\Delta$, $\Theta$, $\Lambda$, $\Sigma$, $\Omega$

\subsection{运算符}

\begin{equation}
  a + b - c \times d \div e = f
\end{equation}

\begin{equation}
  a \pm b \mp c \cdot d / e
\end{equation}

\subsection{关系符号}

\begin{equation}
  a = b \neq c < d \leq e > f \geq g \approx h \equiv i
\end{equation}

\subsection{集合符号}

\begin{equation}
  A \cup B \cap C \subset D \subseteq E \in F \notin G \emptyset
\end{equation}

\subsection{逻辑符号}

\begin{equation}
  \forall x \exists y : x \land y \lor z \neg w \implies v \iff u
\end{equation}

\section{复杂公式}

\subsection{分数和根式}

\begin{equation}
  \frac{a+b}{c+d} = \frac{\sqrt{x^2+y^2}}{\sqrt[3]{z^3+w^3}}
\end{equation}

\subsection{求和与积分}

\begin{equation}
  \sum_{i=1}^{n} i = \frac{n(n+1)}{2}
\end{equation}

\begin{equation}
  \int_0^\infty e^{-x^2} \, dx = \frac{\sqrt{\pi}}{2}
\end{equation}

\begin{equation}
  \prod_{i=1}^{n} x_i = x_1 \cdot x_2 \cdot \ldots \cdot x_n
\end{equation}

\subsection{极限}

\begin{equation}
  \lim_{x \to 0} \frac{\sin x}{x} = 1
\end{equation}

\begin{equation}
  \lim_{n \to \infty} \left(1 + \frac{1}{n}\right)^n = e
\end{equation}

\subsection{矩阵}

使用 \texttt{matrix} 环境：

\begin{equation}
  A = \begin{pmatrix}
    a_{11} & a_{12} & a_{13} \\
    a_{21} & a_{22} & a_{23} \\
    a_{31} & a_{32} & a_{33}
  \end{pmatrix}
\end{equation}

不同的括号类型：

\begin{equation}
  \begin{bmatrix} 1 & 2 \\ 3 & 4 \end{bmatrix}
  \begin{vmatrix} 1 & 2 \\ 3 & 4 \end{vmatrix}
  \begin{Bmatrix} 1 & 2 \\ 3 & 4 \end{Bmatrix}
\end{equation}

\subsection{分段函数}

\begin{equation}
  f(x) = \begin{cases}
    x^2 & \text{if } x \geq 0 \\
    -x^2 & \text{if } x < 0
  \end{cases}
\end{equation}

\section{多行公式}

\subsection{对齐公式}

使用 \texttt{align} 环境：

\begin{align}
  a &= b + c \\
  &= d + e + f \\
  &= g
\end{align}

\subsection{多个公式}

\begin{align}
  x &= a + b \label{eq:x} \\
  y &= c + d \label{eq:y} \\
  z &= e + f \label{eq:z}
\end{align}

可以分别引用：\eqref{eq:x}, \eqref{eq:y}, \eqref{eq:z}。

\subsection{公式组}

使用 \texttt{gather} 环境（居中，不对齐）：

\begin{gather}
  a = b + c \\
  x = y + z \\
  p = q + r
\end{gather}

\section{特殊公式}

\subsection{导数}

\begin{equation}
  \frac{d}{dx}f(x) = f'(x) = \dot{f}
\end{equation}

\begin{equation}
  \frac{\partial f}{\partial x} = f_x
\end{equation}

\subsection{向量}

\begin{equation}
  \vec{v} = \overrightarrow{AB} = (x, y, z)
\end{equation}

\begin{equation}
  \mathbf{v} \cdot \mathbf{w} = |\mathbf{v}| |\mathbf{w}| \cos\theta
\end{equation}

\subsection{概率统计}

\begin{equation}
  P(A|B) = \frac{P(B|A)P(A)}{P(B)}
\end{equation}

\begin{equation}
  \mathbb{E}[X] = \sum_{i=1}^{n} x_i p_i
\end{equation}

\begin{equation}
  \text{Var}(X) = \mathbb{E}[(X - \mathbb{E}[X])^2]
\end{equation}

\subsection{复数}

\begin{equation}
  z = a + bi = r e^{i\theta} = r(\cos\theta + i\sin\theta)
\end{equation}

\section{定理环境}

\begin{theorem}[勾股定理]
  在直角三角形中，斜边的平方等于两直角边的平方和：
  \begin{equation}
    a^2 + b^2 = c^2
  \end{equation}
\end{theorem}

\begin{proof}
  证明过程...
\end{proof}

\begin{lemma}
  引理内容...
\end{lemma}

\begin{corollary}
  推论内容...
\end{corollary}

\begin{definition}
  定义内容...
\end{definition}

\section{常见问题}

\subsection{公式编号}

\begin{itemize}
  \item 使用 \texttt{equation} 环境自动编号
  \item 使用 \texttt{equation*} 或 \texttt{\textbackslash[...\textbackslash]} 不编号
  \item 使用 \texttt{\textbackslash label} 设置标签
  \item 使用 \texttt{\textbackslash ref} 或 \texttt{\textbackslash eqref} 引用
\end{itemize}

\subsection{公式间距}

\begin{itemize}
  \item 使用 \texttt{\textbackslash,} 添加小间距
  \item 使用 \texttt{\textbackslash quad} 添加中等间距
  \item 使用 \texttt{\textbackslash qquad} 添加大间距
\end{itemize}

\subsection{字体}

\begin{itemize}
  \item 普通变量：$x$（斜体）
  \item 向量：$\mathbf{v}$（粗体）
  \item 矩阵：$\mathbf{A}$（粗体）
  \item 集合：$\mathbb{R}$（空心体）
  \item 算子：$\mathrm{d}x$（正体）
\end{itemize}

\section{最佳实践}

\begin{enumerate}
  \item \textbf{使用 equation 环境}：重要公式使用带编号的环境
  \item \textbf{引用公式}：在正文中引用所有重要公式
  \item \textbf{对齐公式}：多行公式使用 align 环境对齐
  \item \textbf{标点符号}：公式后的标点符号放在公式内
  \item \textbf{变量命名}：使用有意义的变量名
  \item \textbf{单位}：使用 \texttt{siunitx} 宏包处理单位
  \item \textbf{间距}：适当使用间距命令改善可读性
\end{enumerate}

\section{推荐宏包}

\begin{itemize}
  \item \texttt{amsmath} - 基础数学宏包（必需）
  \item \texttt{amssymb} - 额外数学符号
  \item \texttt{mathtools} - amsmath 的增强版
  \item \texttt{unicode-math} - Unicode 数学字体
  \item \texttt{siunitx} - 单位和数字格式化
  \item \texttt{physics} - 物理公式
  \item \texttt{bm} - 粗体数学符号
\end{itemize}
